\section{Transporte de Sedimentos}

\subsection{Definição}

O \textbf{transporte de sedimentos} corresponde ao deslocamento de partículas,
fragmentos de rocha, minerais e materiais biogênicos desde a área-fonte até os
locais onde serão posteriormente depositados. O transporte ocorre depois que o
material é produzido ou disponibilizado pelo intemperismo e pela erosão e
constitui uma etapa fundamental do ciclo sedimentar \parencite{nichols2009}.

De maneira simplificada:

\[
\boxed{
\text{intemperismo}
\rightarrow
\text{erosão}
\rightarrow
\text{transporte}
\rightarrow
\text{deposição}
\rightarrow
\text{sedimentação}
}
\]

Durante o transporte, os sedimentos podem sofrer modificações importantes, como:

\begin{itemize}
    \item fragmentação e abrasão;
    \item arredondamento dos grãos;
    \item seleção granulométrica;
    \item separação de minerais em função de densidade, tamanho e resistência;
    \item retrabalhamento sucessivo.
\end{itemize}

Assim, as características de um depósito sedimentar refletem não apenas a
natureza da rocha-fonte, mas também os agentes, a distância, a duração e a
energia envolvida no transporte.

\subsection{Principais agentes de transporte}

Os principais agentes capazes de transportar sedimentos são:

\begin{enumerate}
    \item água;
    \item vento;
    \item gelo;
    \item gravidade.
\end{enumerate}

A água constitui um dos principais agentes de transporte, especialmente nos
sistemas fluviais, lacustres, estuarinos, deltaicos, costeiros e marinhos.

\subsection{Modos de transporte sedimentar}

Os principais modos de transporte são:

\begin{itemize}
    \item tração;
    \item saltação;
    \item suspensão;
    \item transporte em solução.
\end{itemize}

Um mesmo grão pode mudar de modo de transporte quando variam a velocidade, a
turbulência e a energia do fluxo.

\subsubsection{Tração}

A \textbf{tração} corresponde ao transporte de partículas junto ao fundo. Os
grãos podem deslocar-se por rolamento, deslizamento ou arraste.

Esse mecanismo é importante para partículas relativamente grossas ou pesadas,
como areia grossa, grânulos, seixos e cascalhos. Para que o movimento se inicie,
a força exercida pelo fluido deve superar a resistência associada ao peso do
grão, ao atrito com o fundo e ao encaixe entre as partículas.

\subsubsection{Saltação}

Na \textbf{saltação}, as partículas deslocam-se por uma sequência de pequenos
saltos sobre o substrato. O fluxo eleva temporariamente o grão, que percorre uma
determinada distância e retorna ao fundo.

De maneira esquemática:

\[
\text{repouso}
\rightarrow
\text{elevação}
\rightarrow
\text{deslocamento}
\rightarrow
\text{impacto}
\rightarrow
\text{novo salto}
\]

A saltação é especialmente importante no transporte de areia e ocorre tanto em
ambientes subaquosos quanto eólicos. Esse mecanismo participa da formação e da
migração de marcas onduladas e dunas.

\subsubsection{Suspensão}

Na \textbf{suspensão}, as partículas permanecem temporariamente mantidas dentro
da coluna de fluido, principalmente pela turbulência.

Esse mecanismo é especialmente importante para materiais finos, como:

\begin{itemize}
    \item argila;
    \item silte;
    \item areia muito fina.
\end{itemize}

Quanto menor a velocidade de queda da partícula, maior sua tendência a
permanecer em suspensão. Por essa razão, sedimentos muito finos podem ser
transportados por grandes distâncias.

\subsubsection{Transporte em solução}

Parte dos produtos do intemperismo químico é transportada na forma de íons
dissolvidos. Entre os principais componentes transportados em solução estão:

\[
Ca^{2+}, \quad Na^+, \quad K^+, \quad Mg^{2+}
\]

e

\[
HCO_3^-, \quad Cl^-, \quad SO_4^{2-}.
\]

Alterações posteriores das condições físico-químicas da água podem provocar a
precipitação desses componentes, contribuindo para a formação de sedimentos
químicos.

\subsection{Carga sedimentar}

Em sistemas fluviais, a carga transportada pode ser dividida em três componentes
principais.

\subsubsection{Carga de fundo}

A \textbf{carga de fundo} é formada pelo material que permanece próximo ou em
contato frequente com o leito e é transportado principalmente por tração e
saltação.

\subsubsection{Carga em suspensão}

A \textbf{carga em suspensão} corresponde às partículas mantidas na coluna
d'água pela turbulência, principalmente silte, argila e parte da areia fina.

\subsubsection{Carga dissolvida}

A \textbf{carga dissolvida} é constituída por íons e moléculas transportados em
solução.

\subsection{Competência e capacidade de transporte}

Dois conceitos importantes para compreender o comportamento de um fluxo são
\textbf{competência} e \textbf{capacidade}.

A \textbf{competência} corresponde ao maior tamanho de partícula que um fluxo
consegue transportar. Assim, de maneira geral:

\[
\text{maior energia do fluxo}
\Rightarrow
\text{maior competência}
\]

A \textbf{capacidade}, por sua vez, corresponde à quantidade total de sedimento
que o fluxo consegue transportar.

Portanto:

\begin{itemize}
    \item competência relaciona-se principalmente ao tamanho máximo dos grãos;
    \item capacidade relaciona-se principalmente à quantidade total de sedimento.
\end{itemize}

\subsection{Velocidade do fluxo e granulometria}

A erosão, o transporte e a deposição dependem da interação entre:

\begin{itemize}
    \item velocidade do fluxo;
    \item turbulência;
    \item tamanho e densidade dos grãos;
    \item forma das partículas;
    \item viscosidade e densidade do fluido;
    \item coesão do sedimento;
    \item características do fundo.
\end{itemize}

Correntes mais energéticas podem transportar partículas maiores. Entretanto, a
relação entre tamanho do grão e facilidade de erosão não é linear. Partículas
argilosas, apesar de muito pequenas, podem exigir velocidades relativamente
elevadas para serem erodidas devido às forças de coesão entre os grãos.

Assim:

\[
\text{grão pequeno}
\not\Rightarrow
\text{necessariamente fácil erosão}.
\]

\subsection{Número de Froude e regimes de fluxo}

Em canais abertos, o número de Froude compara a velocidade média do fluxo com a velocidade de propagação de ondas gravitacionais rasas:

\[
\mathrm{Fr}=\frac{U}{\sqrt{gD}},
\]

em que \(U\) é a velocidade média, \(g\) é a aceleração da gravidade e \(D\) é uma profundidade característica do fluxo.

\begin{itemize}
  \item \(\mathrm{Fr}<1\): fluxo \textbf{subcrítico}, no qual perturbações podem se propagar para montante;
  \item \(\mathrm{Fr}=1\): condição \textbf{crítica};
  \item \(\mathrm{Fr}>1\): fluxo \textbf{supercrítico}, no qual a velocidade do fluxo supera a das ondas rasas.
\end{itemize}

Em leitos arenosos, o regime inferior inclui comumente ripples, dunas e leito plano inferior. Com o aumento da velocidade e do transporte em suspensão, podem surgir leito plano superior, antidunas e outras formas relacionadas ao regime superior. A forma efetivamente produzida depende também da granulometria, da profundidade e da carga sedimentar \parencite{harms1975,nichols2009}.

\subsection{Diagrama de Hjulström}

A relação entre tamanho das partículas e velocidade da corrente é
tradicionalmente representada pelo \textbf{diagrama de Hjulström}. O diagrama
distingue três campos principais:

\begin{enumerate}
    \item erosão;
    \item transporte;
    \item deposição.
\end{enumerate}

As areias apresentam valores relativamente baixos de velocidade crítica para o
início do movimento. As argilas, apesar de muito finas, podem exigir velocidades
mais elevadas para serem erodidas devido à coesão. Uma vez colocadas em
suspensão, entretanto, podem continuar sendo transportadas mesmo quando a
velocidade do fluxo diminui.

\subsection{Velocidade de queda}

A \textbf{velocidade de queda} corresponde à velocidade com que uma partícula
tende a se deslocar em direção ao fundo sob a ação da gravidade.

Ela depende de fatores como:

\begin{itemize}
    \item tamanho;
    \item densidade;
    \item forma da partícula;
    \item densidade do fluido;
    \item viscosidade do fluido.
\end{itemize}

De forma geral:

\[
\text{partículas maiores e mais densas}
\Rightarrow
\text{maior velocidade de queda}.
\]

Quando a energia de um fluxo diminui, as partículas maiores normalmente tendem
a ser depositadas antes das menores.

\subsection{Seleção granulométrica}

Durante o transporte pode ocorrer \textbf{seleção granulométrica}, pois partículas
de diferentes tamanhos e densidades respondem de maneiras distintas ao mesmo
fluxo.

Um sedimento é considerado:

\begin{itemize}
    \item \textbf{bem selecionado}, quando apresenta pequena variação granulométrica;
    \item \textbf{mal selecionado}, quando apresenta ampla variação dos tamanhos dos grãos.
\end{itemize}

O vento e correntes persistentes podem produzir depósitos relativamente bem
selecionados. Geleiras e determinados fluxos gravitacionais, por outro lado,
podem produzir depósitos muito mal selecionados.

\subsection{Abrasão, arredondamento e maturidade}

Durante o transporte, os grãos podem colidir entre si, com o fundo e com
obstáculos presentes no percurso. Essas colisões provocam \textbf{abrasão}, com
desgaste progressivo de cantos e arestas.

De maneira geral:

\[
\text{maior retrabalhamento}
\Rightarrow
\text{maior tendência ao arredondamento}.
\]

O grau de arredondamento também depende da mineralogia, da dureza, do tamanho do
grão e do mecanismo de transporte.

O retrabalhamento sucessivo pode aumentar a maturidade do sedimento. A
\textbf{maturidade textural} relaciona-se principalmente ao grau de seleção, ao
arredondamento dos grãos e à quantidade de matriz fina. A
\textbf{maturidade mineralógica} relaciona-se à proporção de minerais
resistentes. Minerais pouco estáveis tendem a ser alterados ou destruídos,
enquanto minerais resistentes, como o quartzo, podem tornar-se
proporcionalmente mais abundantes.

\subsection{Transporte pelo vento}

O vento atua principalmente sobre sedimentos secos e pouco coesivos. O
transporte eólico ocorre sobretudo por saltação, suspensão e arraste superficial.

Como o ar possui densidade muito menor que a água, sua competência para
transportar partículas grossas é relativamente limitada. Os depósitos eólicos
são geralmente dominados por areia e silte e podem apresentar boa seleção
granulométrica.

\subsection{Transporte pelo gelo}

As geleiras são capazes de transportar sedimentos com ampla variação
granulométrica, desde argila e silte até grandes blocos de rocha.

O transporte glacial apresenta baixa eficiência de seleção, podendo produzir
depósitos constituídos por partículas de tamanhos muito diferentes.

\subsection{Transporte por gravidade}

A gravidade pode produzir deslocamento de sedimentos ao longo das encostas por
processos como:

\begin{itemize}
    \item queda de blocos;
    \item escorregamentos;
    \item deslizamentos;
    \item fluxos de detritos;
    \item fluxos de lama.
\end{itemize}

Os depósitos resultantes frequentemente apresentam baixa seleção e podem conter
clastos relativamente angulosos.

\subsection{Fluxos gravitacionais subaquosos}

Em ambientes submarinos também podem ocorrer movimentos gravitacionais de
sedimentos. Entre os mais importantes estão as \textbf{correntes de turbidez},
formadas por misturas de água e sedimentos com densidade maior que a água
circundante.

À medida que o fluxo perde energia, as partículas maiores tendem a se depositar
antes das menores, podendo produzir uma gradação normal:

\[
\text{base mais grossa}
\rightarrow
\text{topo mais fino}.
\]

Os depósitos produzidos por correntes de turbidez são chamados de
\textbf{turbiditos}.

\subsection{Transporte em ambientes costeiros}

Nos ambientes costeiros, o transporte sedimentar resulta da interação entre:

\begin{itemize}
    \item ondas;
    \item correntes;
    \item marés;
    \item descargas fluviais;
    \item vento.
\end{itemize}

As ondas podem colocar partículas do fundo em movimento ou suspensão, enquanto
as correntes promovem seu deslocamento. O transporte longitudinal ao longo da
costa exerce importante papel na evolução de praias, pontais e demais feições
sedimentares costeiras.

\subsection{Floculação das argilas}

As partículas de argila apresentam dimensões muito pequenas e podem permanecer
em suspensão por longos períodos. Em ambientes estuarinos, a mistura entre água
doce e água salgada pode favorecer a agregação das partículas finas.

Esse processo é denominado \textbf{floculação}. Os agregados, ou flocos,
apresentam velocidade de queda maior do que as partículas individuais,
favorecendo a deposição de sedimentos finos.

\subsection{Relação entre transporte e deposição}

A deposição ocorre quando o agente de transporte deixa de possuir energia
suficiente para manter determinada partícula em movimento.

De maneira simplificada, quando a energia diminui progressivamente, tende a
ocorrer a seguinte sequência de deposição:

\[
\text{cascalho}
\rightarrow
\text{areia}
\rightarrow
\text{silte}
\rightarrow
\text{argila}.
\]

Essa sequência é uma simplificação, pois coesão, floculação, turbulência e
características específicas do fluxo também influenciam a deposição.


% =========================================================
% ESTRUTURAS SEDIMENTARES
% =========================================================
